% Copyright (c) 2026 Geovana Neves. Licensed under CC BY 4.0 (https://creativecommons.org/licenses/by/4.0/). See LICENSE-AND-AUTHORSHIP.md.
%
% shared/info.tex
%
% Identity of the nine guides, 01 to 09: the version they are written for, the
% authorship and licence notice, the citation, and the generic placeholders
% that stand for every site-specific path. Every deck reads this one file, so
% no two decks can state different versions, notices or paths.

% ------------------------------------------------------------ the series ---
% Printed on every cover: "The pyflightstream guides, guide N of 9,
% numbered 1 to 9". The files are named pyfts-guide-NN-<topic>.pdf.
\newcommand{\guideseries}{The pyflightstream guides}
\newcommand{\guidecount}{9}
\newcommand{\guiderange}{1 to 9}

% ----------------------------------------------- the version written for ---
% One statement per guide, printed on its title page. No history of earlier
% releases: the decks describe this version and nothing else.
\newcommand{\pkgname}{pyflightstream}
\newcommand{\pkgversion}{0.35.0}
\newcommand{\guidewritten}{Written for \pkgname{} \pkgversion}

% The solver build the examples name. Not a version of the package.
\newcommand{\fsbuild}{26.123}

% ----------------------------------------- authorship and the licence ------
\newcommand{\pkgauthor}{Geovana Neves}
\newcommand{\pkglicense}{MIT}
\newcommand{\pkgrepo}{https://github.com/nevesgeovana/pyflightstream}
% The concept DOI resolves to the latest archived release of the package.
\newcommand{\pkgdoiconcept}{10.5281/zenodo.21482924}

% The package's documentation site. A deck cites one page of it as a
% numbered reference, by the page's slug (its file name under docs/ without
% .md) and its title:
%   \item \docref{storage-and-sync}{Storage and sync}
% A tier-1 test of the repository checks that every slug names a page that
% exists under docs/.
\newcommand{\pkgsite}{https://nevesgeovana.github.io/pyflightstream}
\newcommand{\docref}[2]{\pkgauthor, \emph{\pkgname{} documentation}, page
  ``#2''. \url{\pkgsite/#1/}}

% The measurement records a deck cites for the solver's measured behaviour.
\newcommand{\recordsref}{\pkgauthor, \emph{\pkgname: measurement records of
  the solver's behaviour}, 2026. \url{\pkgrepo}, concept DOI
  \texttt{\pkgdoiconcept}.}

% The public airframe geometry the decks point to for practice.
\newcommand{\smiref}{G. Neves, ``SMI: airframe geometry'', version 1.0.0,
  Zenodo, 2026. \url{https://zenodo.org/records/23025105}}

\newcommand{\guidenotice}{Copyright (c) 2026 Geovana Neves. Licensed under
  CC BY 4.0 (\url{https://creativecommons.org/licenses/by/4.0/}). See
  LICENSE-AND-AUTHORSHIP.md.}
\newcommand{\guidefootnotice}{Copyright (c) 2026 Geovana Neves, CC BY 4.0}

% One block, used wherever a deck asks somebody to cite the tool.
\newcommand{\pkgcitation}{%
  \pkgauthor. \emph{\pkgname: version-aware, didactic Python driver for the
  FlightStream panel-method solver}, 2026. \url{\pkgrepo}. Concept DOI:
  \texttt{\pkgdoiconcept}.}

% ------------------------------------------- the generic placeholders ------
% Every site-specific path appears in the decks through these macros and
% nowhere as a literal. Each site keeps its own values in a local note; the
% decks stay the same everywhere.
%
%   \mainworkspace       the main workspace on the Windows workstation
%   \clusterroot         the shared install folder on the cluster file system
%   \clusterworkspace    the cluster workspace (runs submitted there)
%   \shareunc            the shared install folder seen from Windows (a network share)
%   \clusterworkspaceunc the cluster workspace seen from Windows, as sync reads it
%   \envname             the name of one versioned environment (venv) folder
\newcommand{\mainworkspace}{C:/<work-folder>/fts-workspace}
\newcommand{\clusterroot}{<cluster-root>/fts-unix}
\newcommand{\clusterworkspace}{<cluster-root>/fts-unix/fts-workspace}
\newcommand{\shareunc}{\ftsbs\ftsbs<server>\ftsbs<share>\ftsbs fts-unix}
\newcommand{\clusterworkspaceunc}{//<server>/<share>/fts-unix/fts-workspace}
\newcommand{\envname}{pfs0340}
